% !TeX encoding = UTF-8
% !TeX program = xelatex
% !TeX spellcheck = en_US

\documentclass[degree=doctor]{thuthesis}
  % 学位 degree:
  %   doctor | master | bachelor | postdoc
  % 学位类型 degree-type:
  %   academic（默认）| professional
  % 语言 language
  %   chinese（默认）| english
  % 字体库 fontset
  %   windows | mac | fandol | ubuntu
  % 建议终版使用 Windows 平台的字体编译

% 论文基本配置，加载宏包等全局配置

\usepackage{amsmath}
\numberwithin{equation}{section}
\usepackage{array}
\usepackage{caption}
\usepackage{graphicx}
\usepackage{siunitx}
\usepackage{colortbl}
\usepackage{multirow}
\usepackage{hhline}
\usepackage{calc}
\usepackage{tabularx}
\usepackage{threeparttable}
\usepackage{wrapfig}
\renewcommand{\topfraction}{.85}
\renewcommand{\bottomfraction}{.7}
\renewcommand{\textfraction}{.15}
\renewcommand{\floatpagefraction}{.66}
\usepackage{pdflscape}
\usepackage{lscape}
\newcommand{\blandscape}{\begin{landscape}}
\newcommand{\elandscape}{\end{landscape}}
\setcounter{topnumber}{3}
\setcounter{bottomnumber}{3}
\setcounter{totalnumber}{4}
\usepackage{ctex}
\newlength{\Oldarrayrulewidth}
\setlength{\Oldarrayrulewidth}{\arrayrulewidth}
\newlength{\Oldtabcolsep}
\setlength{\Oldtabcolsep}{\tabcolsep}
\usepackage{bookmark}

% !TeX root = ./thuthesis-example.tex

% 论文基本信息配置

\thusetup{
  %******************************
  % 注意：
  %   1. 配置里面不要出现空行
  %   2. 不需要的配置信息可以删除
  %   3. 建议先阅读文档中所有关于选项的说明
  %******************************
  %
  % 输出格式
  %   选择打印版（print）或用于提交的电子版（electronic），前者会插入空白页以便直接双面打印
  %
  output = print,
  %
  % 标题
  %   可使用“\\”命令手动控制换行
  %
  title  = {清华大学学位论文 Quarto 模板\\使用示例文档},
  title* = {An Introduction to Quarto Thesis Template of Tsinghua
            University},
  %
  % 学科门类
  %   1. 学术型
  %      - 中文
  %        需注明所属的学科门类，例如：
  %        哲学、经济学、法学、教育学、文学、历史学、理学、工学、农学、医学、
  %        军事学、管理学、艺术学
  %      - 英文
  %        博士：Doctor of Philosophy
  %        硕士：
  %          哲学、文学、历史学、法学、教育学、艺术学门类，公共管理学科
  %          填写“Master of Arts“，其它填写“Master of Science”
  %   2. 专业型
  %      直接填写专业学位的名称，例如：
  %      教育博士、工程硕士等
  %      Doctor of Education, Master of Engineering
  %   3. 本科生不需要填写
  %
  degree-category  = {法学博士},
  degree-category* = {Doctor of Philosophy},
  %
  % 培养单位
  %   填写所属院系的全名
  %
  department = {社会科学学院},
  %
  % 学科
  %   1. 研究生学术型学位，获得一级学科授权的学科填写一级学科名称，其他填写二级学科名称
  %   2. 本科生填写专业名称，第二学位论文需标注“（第二学位）”
  %
  discipline  = {政治学},
  discipline* = {Political Science},
  %
  % 专业领域
  %   1. 设置专业领域的专业学位类别，填写相应专业领域名称
  %   2. 2019 级及之前工程硕士学位论文，在 `engineering-field` 填写相应工程领域名称
  %   3. 其他专业学位类别的学位论文无需此信息
  %
  % professional-field  = {计算机技术},
  % professional-field* = {Computer Technology},
  %
  % 姓名
  %
  author  = {孙宇飞},
  author* = {Sun Yufei},
  %
  % 指导教师
  %   中文姓名和职称之间以英文逗号“,”分开，下同
  %
  supervisor  = {胡悦, 副教授},
  supervisor* = {Associate Professor Hu Yue},
  %
  % 副指导教师
  %
  % associate-supervisor  = {XXX, 教授},
  % associate-supervisor* = {Professor XXX},
  %
  % 联合指导教师
  %
  % co-supervisor  = {某某某, 教授},
  % co-supervisor* = {Professor Mou Moumou},
  %
  % 日期
  %   使用 ISO 格式；默认为当前时间
  %
  % date = {2019-07-07},
  %
  % 是否在中文封面后的空白页生成书脊（默认 false）
  %
  include-spine = false,
  %
  % 密级和年限
  %   秘密, 机密, 绝密
  %
  % secret-level = {秘密},
  % secret-year  = {10},
  %
  % 博士后专有部分
  %
  % clc                = {分类号},
  % udc                = {UDC},
  % id                 = {编号},
  % discipline-level-1 = {计算机科学与技术},  % 流动站（一级学科）名称
  % discipline-level-2 = {系统结构},          % 专业（二级学科）名称
  % start-date         = {2011-07-01},        % 研究工作起始时间
}

% 载入所需的宏包

% 定理类环境宏包
\usepackage{amsthm}
% 也可以使用 ntheorem
% \usepackage[amsmath,thmmarks,hyperref]{ntheorem}

\thusetup{
  %
  % 数学字体
  % math-style = GB,  % GB | ISO | TeX
  math-font  = xits,  % stix | xits | libertinus
}

% 可以使用 nomencl 生成符号和缩略语说明
% \usepackage{nomencl}
% \makenomenclature

% 表格加脚注
\usepackage{threeparttable}

% 表格中支持跨行
\usepackage{multirow}
\providecommand{\tightlist}{%
  \setlength{\itemsep}{0pt}\setlength{\parskip}{0pt}}

% 固定宽度的表格。
% \usepackage{tabularx}

% 跨页表格
\usepackage{longtable}

% 算法
\usepackage{algorithm}
\usepackage{algorithmic}

% 量和单位
\usepackage{siunitx}

% 参考文献使用 BibTeX + natbib 宏包
% 顺序编码制
\usepackage[sort]{natbib}
\bibliographystyle{thuthesis-numeric}

% 著者-出版年制
% \usepackage{natbib}
% \bibliographystyle{thuthesis-author-year}

% 本科生参考文献的著录格式
% \usepackage[sort]{natbib}
% \bibliographystyle{thuthesis-bachelor}

% 参考文献使用 BibLaTeX 宏包
% \usepackage[style=thuthesis-numeric]{biblatex}
% \usepackage[style=thuthesis-author-year]{biblatex}
% \usepackage[style=gb7714-2015]{biblatex}
% \usepackage[style=apa]{biblatex}
% \usepackage[style=mla-new]{biblatex}
% 声明 BibLaTeX 的数据库
% \addbibresource{ref/refs.bib}

% 定义所有的图片文件在 figures 子目录下
\graphicspath{{figures/}}

% 数学命令
\makeatletter
\newcommand\dif{%  % 微分符号
  \mathop{}\!%
  \ifthu@math@style@TeX
    d%
  \else
    \mathrm{d}%
  \fi
}
\makeatother

% hyperref 宏包在最后调用
\usepackage{hyperref}


\begin{document}

% 封面
\maketitle

% 学位论文指导小组、公开评阅人和答辩委员会名单
% 本科生不需要
% !TeX root = ../thuthesis-example.tex

\begin{committee}[name={学位论文指导小组、公开评阅人和答辩委员会名单}]

  \newcolumntype{C}[1]{@{}>{\centering\arraybackslash}p{#1}}

  \section*{指导小组名单}

  \begin{center}
    \begin{tabular}{C{3cm}C{3cm}C{9cm}@{}}
      李XX & 教授     & 清华大学 \\
      王XX & 副教授   & 清华大学 \\
      张XX & 助理教授 & 清华大学 \\
    \end{tabular}
  \end{center}


  \section*{公开评阅人名单}

  \begin{center}
    \begin{tabular}{C{3cm}C{3cm}C{9cm}@{}}
      刘XX & 教授   & 清华大学                    \\
      陈XX & 副教授 & XXXX大学                    \\
      杨XX & 研究员 & 中国XXXX科学院XXXXXXX研究所 \\
    \end{tabular}
  \end{center}


  \section*{答辩委员会名单}

  \begin{center}
    \begin{tabular}{C{2.75cm}C{2.98cm}C{4.63cm}C{4.63cm}@{}}
      主席 & 赵XX                  & 教授                    & 清华大学       \\
      委员 & 刘XX                  & 教授                    & 清华大学       \\
          & \multirow{2}{*}{杨XX} & \multirow{2}{*}{研究员} & 中国XXXX科学院 \\
          &                       &                         & XXXXXXX研究所  \\
          & 黄XX                  & 教授                    & XXXX大学       \\
          & 周XX                  & 副教授                  & XXXX大学       \\
      秘书 & 吴XX                  & 助理研究员              & 清华大学       \\
    \end{tabular}
  \end{center}

\end{committee}



% 也可以导入 Word 版转的 PDF 文件
% \begin{committee}[file=figures/committee.pdf]
% \end{committee}


% 使用授权的说明
\copyrightpage
% 将签字扫描后授权文件 scan-copyright.pdf 替换原始页面
% \copyrightpage[file=scan-copyright.pdf]

\frontmatter
% !TeX root = ../thuthesis-example.tex

% 中英文摘要和关键字

\begin{abstract}
  论文的摘要是对论文研究内容和成果的高度概括。
  摘要应对论文所研究的问题及其研究目的进行描述，对研究方法和过程进行简单介绍，对研究成果和所得结论进行概括。
  摘要应具有独立性和自明性，其内容应包含与论文全文同等量的主要信息。
  使读者即使不阅读全文，通过摘要就能了解论文的总体内容和主要成果。

  论文摘要的书写应力求精确、简明。
  切忌写成对论文书写内容进行提要的形式，尤其要避免“第 1 章……；第 2 章……；……”这种或类似的陈述方式。

  关键词是为了文献标引工作、用以表示全文主要内容信息的单词或术语。
  关键词不超过 5 个，每个关键词中间用分号分隔。

  % 关键词用“英文逗号”分隔，输出时会自动处理为正确的分隔符
  \thusetup{
    keywords = {关键词 1, 关键词 2, 关键词 3, 关键词 4, 关键词 5},
  }
\end{abstract}

\begin{abstract*}
  An abstract of a dissertation is a summary and extraction of research work and contributions.
  Included in an abstract should be description of research topic and research objective, brief introduction to methodology and research process, and summary of conclusion and contributions of the research.
  An abstract should be characterized by independence and clarity and carry identical information with the dissertation.
  It should be such that the general idea and major contributions of the dissertation are conveyed without reading the dissertation.

  An abstract should be concise and to the point.
  It is a misunderstanding to make an abstract an outline of the dissertation and words “the first chapter”, “the second chapter” and the like should be avoided in the abstract.

  Keywords are terms used in a dissertation for indexing, reflecting core information of the dissertation.
  An abstract may contain a maximum of 5 keywords, with semi-colons used in between to separate one another.

  % Use comma as separator when inputting
  \thusetup{
    keywords* = {keyword 1, keyword 2, keyword 3, keyword 4, keyword 5},
  }
\end{abstract*}


% 目录
\tableofcontents

% 插图和附表清单
% 本科生的插图索引和表格索引需要移至正文之后、参考文献前
% \listoffiguresandtables  % 插图和附表清单（仅限研究生）
\listoffigures           % 插图清单
\listoftables            % 附表清单

% 符号对照表
% !TeX root = ../thuthesis-example.tex

\begin{denotation}[3cm]
  \item[CCP] 中国共产党（Chinese Communist Party）
  \item[CPC] 中国共产党（Chinese Communist Party）[注：CCP和CPC都可以指中国共产党]
  \item[NPC] 全国人民代表大会（National People's Congress）
  \item[CPPCC] 全国政协（Chinese People's Political Consultative Conference）
  \item[PLA] 中国人民解放军（People's Liberation Army）
  \item[CASS] 中国社会科学院（Chinese Academy of Social Sciences）
  \item[CISA] 中国国际战略学会（China Institute of International Strategic Studies）
  \item[RF] 改革派（Reformist Faction）
  \item[CF] 保守派（Conservative Faction）
  \item[FTZ] 自贸区（Free Trade Zone）
  \item[BRI] 一带一路（Belt and Road Initiative）
  \item[SOE] 国有企业（State-owned Enterprises）
  \item[CMI] 中央军委（Central Military Commission）
  \item[PBOC] 中国人民银行（People's Bank of China）
  \item[CAC] 网络安全和信息化委员会（Cyberspace Administration of China）
  \item[SAR] 特别行政区（Special Administrative Region）
  \item[CFA] 中央财经委员会（Central Finance and Economics Commission）
  \item[HKSAR] 香港特别行政区（Hong Kong Special Administrative Region）
  \item[MSAR] 澳门特别行政区（Macao Special Administrative Region）
  \item[TAR] 西藏自治区（Tibet Autonomous Region）
  \item[XUAR] 新疆维吾尔自治区（Xinjiang Uyghur Autonomous Region）
  \item[CACG] 中国反腐败工作领导小组（China Anti-Corruption Coordination Group）
  \item[RMB] 人民币（Renminbi）
\end{denotation}




% 也可以使用 nomencl 宏包，需要在导言区
% \usepackage{nomencl}
% \makenomenclature

% 在这里输出符号说明
% \printnomenclature[3cm]

% 在正文中的任意为都可以标题
% \nomenclature{PI}{聚酰亚胺}
% \nomenclature{MPI}{聚酰亚胺模型化合物，N-苯基邻苯酰亚胺}
% \nomenclature{PBI}{聚苯并咪唑}
% \nomenclature{MPBI}{聚苯并咪唑模型化合物，N-苯基苯并咪唑}
% \nomenclature{PY}{聚吡咙}
% \nomenclature{PMDA-BDA}{均苯四酸二酐与联苯四胺合成的聚吡咙薄膜}
% \nomenclature{MPY}{聚吡咙模型化合物}
% \nomenclature{As-PPT}{聚苯基不对称三嗪}
% \nomenclature{MAsPPT}{聚苯基不对称三嗪单模型化合物，3,5,6-三苯基-1,2,4-三嗪}
% \nomenclature{DMAsPPT}{聚苯基不对称三嗪双模型化合物（水解实验模型化合物）}
% \nomenclature{S-PPT}{聚苯基对称三嗪}
% \nomenclature{MSPPT}{聚苯基对称三嗪模型化合物，2,4,6-三苯基-1,3,5-三嗪}
% \nomenclature{PPQ}{聚苯基喹噁啉}
% \nomenclature{MPPQ}{聚苯基喹噁啉模型化合物，3,4-二苯基苯并二嗪}
% \nomenclature{HMPI}{聚酰亚胺模型化合物的质子化产物}
% \nomenclature{HMPY}{聚吡咙模型化合物的质子化产物}
% \nomenclature{HMPBI}{聚苯并咪唑模型化合物的质子化产物}
% \nomenclature{HMAsPPT}{聚苯基不对称三嗪模型化合物的质子化产物}
% \nomenclature{HMSPPT}{聚苯基对称三嗪模型化合物的质子化产物}
% \nomenclature{HMPPQ}{聚苯基喹噁啉模型化合物的质子化产物}
% \nomenclature{PDT}{热分解温度}
% \nomenclature{HPLC}{高效液相色谱（High Performance Liquid Chromatography）}
% \nomenclature{HPCE}{高效毛细管电泳色谱（High Performance Capillary lectrophoresis）}
% \nomenclature{LC-MS}{液相色谱-质谱联用（Liquid chromatography-Mass Spectrum）}
% \nomenclature{TIC}{总离子浓度（Total Ion Content）}
% \nomenclature{\textit{ab initio}}{基于第一原理的量子化学计算方法，常称从头算法}
% \nomenclature{DFT}{密度泛函理论（Density Functional Theory）}
% \nomenclature{$E_a$}{化学反应的活化能（Activation Energy）}
% \nomenclature{ZPE}{零点振动能（Zero Vibration Energy）}
% \nomenclature{PES}{势能面（Potential Energy Surface）}
% \nomenclature{TS}{过渡态（Transition State）}
% \nomenclature{TST}{过渡态理论（Transition State Theory）}
% \nomenclature{$\increment G^\neq$}{活化自由能（Activation Free Energy）}
% \nomenclature{$\kappa$}{传输系数（Transmission Coefficient）}
% \nomenclature{IRC}{内禀反应坐标（Intrinsic Reaction Coordinates）}
% \nomenclature{$\nu_i$}{虚频（Imaginary Frequency）}
% \nomenclature{ONIOM}{分层算法（Our own N-layered Integrated molecular Orbital and molecular Mechanics）}
% \nomenclature{SCF}{自洽场（Self-Consistent Field）}
% \nomenclature{SCRF}{自洽反应场（Self-Consistent Reaction Field）}



% 正文部分
\mainmatter

\hypertarget{ux5f15ux8a00}{%
\chapter{引言}\label{ux5f15ux8a00}}

Scroll down for the English version of README.

本文是R语言软件包\texttt{THUDown}的使用手册，从架构、安装方法和使用方法三个方面对\texttt{THUDown}进行了介绍。\footnote{此文撰写的得到了清华大学赵赫和成功两位同学的支持，其中赵赫对``使用方法''章节的内容提供了主要的支持，成功同学对``安装方法''章节提供了支持。}

\texttt{THUDown}是\texttt{Tsinghua\ University}和\texttt{markdown}的缩写，它是一个基于
\LaTeX 和
Quarto开发的清华大学学位论文Quarto模板，包括本科综合论文训练、硕士论文、博士论文以及博士后出站报告。它旨在让用户能够使用语法简单的类markdown文档（Quarto）来实现调用
\LaTeX 进行复杂的文件编译。\texttt{THUDown}亦被封装在\href{https://sammo3182.github.io/software/drhutools/}{\texttt{drhutools}}软件包中，用户可借助\texttt{drhutools::THUDown()}进行下载使用。

\textbf{需要特别强调的是}，本软件包是基于\href{https://github.com/tuna/thuthesis}{ThuThesis}开发的Quarto版本，因此其亦可被视作其拓展包。
\texttt{THUDown}的\textbf{主要}
\LaTeX 代码来自于\textbf{ThuThesis}宏包开发者的贡献，\texttt{THUDown}在\href{https://www.latex-project.org/lppl/lppl-1-3c/}{LaTeX项目公共许可证
v1.3c}的授权使用范围内使用\textbf{ThuThesis}宏包的内容。

请在\href{https://www.latex-project.org/lppl/lppl-1-3c/}{LaTeX项目公共许可证
v1.3c}的使用许可范围内使用\texttt{THUDown}。

\hypertarget{ux67b6ux6784}{%
\section{架构}\label{ux67b6ux6784}}

\texttt{THUdown}
的编译框架采用了一个分层和集成的方法来生成学术文档。其核心编译流程可以概括为以下几个阶段：

\begin{enumerate}
\def\labelenumi{\arabic{enumi}.}
\item
  \textbf{内容创建与转换}: 初始内容以 Quarto
  文件格式撰写。这种格式主要针对文本的表示和结构化，使其便于编辑和版本控制。
\item
  \textbf{中间表示}: 利用 \texttt{pandoc} 这一强大的文档转换工具，Quarto
  文件被转化为 \LaTeX 格式。在这一阶段，文本和元数据被转化为
  \LaTeX 指令和环境，为下一步套用模板做准备。
\item
  \textbf{模板应用}: 转换得到的 \LaTeX 文件被整合到 \texttt{ThuThesis}
  模板中。\texttt{ThuThesis} 是为清华大学学术论文设计的
  \LaTeX 模板，能确保文档满足清华大学的格式要求。
\item
  \textbf{文档编译}: 最后，通过使用 \texttt{make} 命令（该命令引导
  \LaTeX 编译器），将 \LaTeX 文件编译成为最终的 PDF 格式的学术文档。
\end{enumerate}

这一架构确保了从内容创建到最终文档输出的整个过程都是高度自动化的，同时也保证了文档质量和格式的一致性。


\hypertarget{ux4e0bux8f7d}{%
\chapter{下载}\label{ux4e0bux8f7d}}

\texttt{THUDown}是一个基于 \LaTeX 和
Quarto开发的清华大学学位论文Quarto模板，因此其使用需要下载\texttt{THUDown}核心软件包，亦需要恰当安装
\LaTeX 、Quarto和ThuThesis等支持项目

\hypertarget{ux6838ux5fc3ux8f6fux4ef6ux5305}{%
\section{核心软件包}\label{ux6838ux5fc3ux8f6fux4ef6ux5305}}

本节提供了\textbf{开发版}和\textbf{发布版}下载路径。

\begin{itemize}
\item
  开发版

  \begin{itemize}
  \tightlist
  \item
    \href{https://github.com/syfyufei}{GitHub Releases}
  \end{itemize}
\item
  发布版

  \begin{itemize}
  \item
    \href{}{CRAN}
  \item
    \href{https://sammo3182.github.io/software/drhutools/}{\texttt{drhutools}}
  \end{itemize}
\end{itemize}

\hypertarget{ux652fux6301ux9879ux76ee}{%
\section{支持项目}\label{ux652fux6301ux9879ux76ee}}

\hypertarget{ux53d1ux884cux7248}{%
\subsection{\texorpdfstring{\LaTeX 发行版}{发行版}}\label{ux53d1ux884cux7248}}

本文提供两种不同的 \LaTeX 发行版选项，用户可以根据自身需求选择安装

\begin{itemize}
\item
  \textbf{\texttt{TinyTeX}}。TinyTeX是一个瘦身版的TeX
  Live，用户可以参考其主页的\href{https://yihui.org/tinytex/}{支持文档}自行安装。需要强调的是，因为\texttt{THUDown}和ThuThesis中包含大量宏包的调用，因此安装\texttt{TinyTeX}的用户在享受其便捷的同时，亦需要在安装第三方宏包上伤脑筋，尤其是无法通过\texttt{tlmgr}安装的宏包。对于这些包，请参见\href{https://github.com/rstudio/tinytex/issues/126\#issuecomment-503020154}{此issue}
\item
  \textbf{\texttt{TexLive}}。用户亦可通过\href{https://tug.org/texlive/}{Texlive主页}安装完整版的\texttt{TexLive}。
\end{itemize}

\hypertarget{quarto}{%
\subsection{Quarto}\label{quarto}}

Quarto是一套开源的科技出版系统。
它使得用户能够使用Jupyter、Python、R、Julia和Observable创作HTML、PDF、MS
Word、ePub等格式的可复制、高质量的文章、演示文稿、网站、博客和书籍。
Quarto是\texttt{THUDown}必须项，请参考\href{https://quarto.org/docs/get-started/}{Quarto支持页面}进行安装。

\hypertarget{thuthesis}{%
\subsection{ThuThesis}\label{thuthesis}}

\href{https://github.com/tuna/thuthesis}{ThuThesis}一个简单易用的清华大学学位论文
LaTeX 模板，\texttt{THUDown}的主体
\LaTeX 代码来自ThuThesis。请参考\href{https://github.com/tuna/thuthesis\#readme}{ThuThesis的Github项目主页}进行安装。


\hypertarget{ux4f7fux7528ux8bf4ux660e}{%
\chapter{使用说明}\label{ux4f7fux7528ux8bf4ux660e}}

本章将从编辑和编译两个部分对\texttt{THUdown}的使用方法进行详细的说明。

\hypertarget{ux7f16ux8f91}{%
\section{编辑}\label{ux7f16ux8f91}}

\hypertarget{ux4e3bux4f53ux67b6ux6784}{%
\subsection{主体架构}\label{ux4e3bux4f53ux67b6ux6784}}

├── abstract.tex

├── acknowledgements.tex

├── appendix-survey.tex

├── appendix-translation.tex

├── appendix.tex

├── comments.tex

├── committee.tex

├── denotation.tex

├── main.tex

├── /quarto\_file

├── resolution.tex

└── resume.tex

\hypertarget{ux5bf9quarto.qmdux6587ux6863ux7684ux7f16ux8f91}{%
\subsection{对Quarto(.qmd)文档的编辑}\label{ux5bf9quarto.qmdux6587ux6863ux7684ux7f16ux8f91}}

data/quarto\_file文件夹

正文部分

具体功能的使用

附录部分

\hypertarget{ux5bf9latex.texux6587ux6863ux7684ux7f16ux8f91}{%
\subsection{对Latex(.tex)文档的编辑}\label{ux5bf9latex.texux6587ux6863ux7684ux7f16ux8f91}}

data文件夹中的其他文件

\hypertarget{ux7f16ux8bd1}{%
\section{编译}\label{ux7f16ux8bd1}}

\hypertarget{ux7f16ux8bd1ux65b9ux6cd5}{%
\subsection{编译方法}\label{ux7f16ux8bd1ux65b9ux6cd5}}

\begin{enumerate}
\def\labelenumi{\arabic{enumi}.}
\item
  确保你的系统已经安装了\texttt{zsh}，\texttt{quarto}，\texttt{awk}，和\texttt{make}这些工具。
\item
  打开终端，导航（\texttt{cd}）到\texttt{knit.sh}文件所在的目录，即\texttt{THUdown}的根目录。
\item
  运行以下命令使\texttt{knit.sh}文件变为可执行：
\end{enumerate}

\texttt{\{bash\ eval=FALSE,\ include=FALSE\}\ chmod\ +x\ knit.sh}

\begin{enumerate}
\def\labelenumi{\arabic{enumi}.}
\setcounter{enumi}{3}
\tightlist
\item
  执行脚本来编译你的\texttt{THUdown}包：
\end{enumerate}

\texttt{\{bash\ eval=FALSE,\ include=FALSE\}\ ./knit.sh}

\begin{enumerate}
\def\labelenumi{\arabic{enumi}.}
\setcounter{enumi}{4}
\tightlist
\item
  脚本会按照上述的编译原理进行操作，并最终生成你的论文。
\end{enumerate}





% 其他部分
\backmatter

% 参考文献
\bibliography{ref/peer.bib}  % 参考文献使用 BibTeX 编译
% \printbibliography       % 参考文献使用 BibLaTeX 编译

% 附录
% 本科生需要将附录放到声明之后，个人简历之前
\appendix
% % !TeX root = ../thuthesis-example.tex

\begin{survey}
\label{cha:survey}

\title{Title of the Survey}
\maketitle


\tableofcontents


本科生的外文资料调研阅读报告。


\section{Figures and Tables}

\subsection{Figures}

An example figure in appendix (Figure~\ref{fig:appendix-survey-figure}).

\begin{figure}
  \centering
  \includegraphics[width=0.6\linewidth]{example-image-a.pdf}
  \caption{Example figure in appendix}
  \label{fig:appendix-survey-figure}
\end{figure}


\subsection{Tables}

An example table in appendix (Table~\ref{tab:appendix-survey-table}).

\begin{table}
  \centering
  \caption{Example table in appendix}
  \begin{tabular}{ll}
    \toprule
    File name       & Description                                         \\
    \midrule
    thuthesis.dtx   & The source file including documentaion and comments \\
    thuthesis.cls   & The template file                                   \\
    thuthesis-*.bst & BibTeX styles                                       \\
    thuthesis-*.bbx & BibLaTeX styles for bibliographies                  \\
    thuthesis-*.cbx & BibLaTeX styles for citations                       \\
    \bottomrule
  \end{tabular}
  \label{tab:appendix-survey-table}
\end{table}


\section{Equations}

An example equation in appendix (Equation~\eqref{eq:appendix-survey-equation}).
\begin{equation}
  \frac{1}{2 \uppi \symup{i}} \int_\gamma f = \sum_{k=1}^m n(\gamma; a_k) \mathscr{R}(f; a_k)
  \label{eq:appendix-survey-equation}
\end{equation}


\section{Citations}

Example\cite{dupont1974bone} citations\cite{merkt1995rotational} in appendix
\cite{dupont1974bone,merkt1995rotational}.


% 默认使用正文的参考文献样式；
% 如果使用 BibTeX，可以切换为其他兼容 natbib 的 BibTeX 样式。
\bibliographystyle{unsrtnat}
% \bibliographystyle{IEEEtranN}

% 默认使用正文的参考文献 .bib 数据库；
% 如果使用 BibTeX，可以改为指定数据库，如 \bibliography{ref/refs}。
\printbibliography

\end{survey}
       % 本科生：外文资料的调研阅读报告
% % !TeX root = ../thuthesis-example.tex

\begin{translation}
\label{cha:translation}

\title{书面翻译题目}
\maketitle

\tableofcontents


本科生的外文资料书面翻译。


\section{图表示例}

\subsection{图}

附录中的图片示例（图~\ref{fig:appendix-translation-figure}）。

\begin{figure}
  \centering
  \includegraphics[width=0.6\linewidth]{example-image-a.pdf}
  \caption{附录中的图片示例}
  \label{fig:appendix-translation-figure}
\end{figure}


\subsection{表格}

附录中的表格示例（表~\ref{tab:appendix-translation-table}）。

\begin{table}
  \centering
  \caption{附录中的表格示例}
  \begin{tabular}{ll}
    \toprule
    文件名          & 描述                         \\
    \midrule
    thuthesis.dtx   & 模板的源文件，包括文档和注释 \\
    thuthesis.cls   & 模板文件                     \\
    thuthesis-*.bst & BibTeX 参考文献表样式文件    \\
    thuthesis-*.bbx & BibLaTeX 参考文献表样式文件  \\
    thuthesis-*.cbx & BibLaTeX 引用样式文件        \\
    \bottomrule
  \end{tabular}
  \label{tab:appendix-translation-table}
\end{table}


\section{数学公式}

附录中的数学公式示例（公式\eqref{eq:appendix-translation-equation}）。
\begin{equation}
  \frac{1}{2 \uppi \symup{i}} \int_\gamma f = \sum_{k=1}^m n(\gamma; a_k) \mathscr{R}(f; a_k)
  \label{eq:appendix-translation-equation}
\end{equation}


\section{文献引用}

附录\cite{dupont1974bone}中的参考文献引用\cite{merkt1995rotational}示例
\cite{dupont1974bone,merkt1995rotational}。


\appendix

\section{附录}

附录的内容。


% 书面翻译的参考文献
% 默认使用正文的参考文献样式；
% 如果使用 BibTeX，可以切换为其他兼容 natbib 的 BibTeX 样式。
\bibliographystyle{unsrtnat}
% \bibliographystyle{IEEEtranN}

% 默认使用正文的参考文献 .bib 数据库；
% 如果使用 BibTeX，可以改为指定数据库，如 \bibliography{ref/refs}。
\printbibliography

% 书面翻译对应的原文索引
\begin{translation-index}
  \nocite{mellinger1996laser}
  \nocite{bixon1996dynamics}
  \nocite{carlson1981two}
  \bibliographystyle{unsrtnat}
  \printbibliography
\end{translation-index}

\end{translation}
  % 本科生：外文资料的书面翻译
\hypertarget{sec-appendix}{%
\chapter{补充内容}\label{sec-appendix}}

附录是与论文内容密切相关、但编入正文又影响整篇论文编排的条理和逻辑性的资料，例如某些重要的数据表格、计算程序、统计表等，是论文主体的补充内容，可根据需要设置。

附录中的图、表、数学表达式、参考文献等另行编序号，与正文分开，一律用阿拉伯数字编码，
但在数码前冠以附录的序号。





% 致谢
% !TeX root = ../thuthesis-example.tex

\begin{acknowledgements}
  衷心感谢导师×××教授和物理系××副教授对本人的精心指导。他们的言传身教将使我终生受益。

  在美国麻省理工学院化学系进行九个月的合作研究期间，承蒙 Robert Field 教授热心指导与帮助，不胜感激。

  感谢×××××实验室主任×××教授，以及实验室全体老师和同窗们学的热情帮助和支持！

  本课题承蒙国家自然科学基金资助，特此致谢。
\end{acknowledgements}


% 声明
\statement
% 将签字扫描后的声明文件 scan-statement.pdf 替换原始页面
% \statement[file=scan-statement.pdf]
% 本科生编译生成的声明页默认不加页脚，插入扫描版时再补上；
% 研究生编译生成时有页眉页脚，插入扫描版时不再重复。
% 也可以手动控制是否加页眉页脚
% \statement[page-style=empty]
% \statement[file=scan-statement.pdf, page-style=plain]

% 个人简历、在学期间完成的相关学术成果
% 本科生可以附个人简历，也可以不附个人简历
% !TeX root = ../thuthesis-example.tex

\begin{resume}

  \section*{个人简历}

  199× 年 ×× 月 ×× 日出生于河南××市。

  2016 年 9 月考入××大学化学系××国际政治专业，2020 年 6 月本科毕业并获得法学学士学位。

  2020 年 9 月免试进入清华大学社科学院政治学系攻读大数据政治学博士至今。


  \section*{在学期间完成的相关学术成果}

  \subsection*{学术论文}

  \begin{achievements}
    \item Yang Y, Ren T L, Zhang L T, et al. Miniature microphone with silicon-based ferroelectric thin films[J]. Integrated Ferroelectrics, 2003, 52:229-235.
    \item 杨轶, 张宁欣, 任天令, 等. 硅基铁电微声学器件中薄膜残余应力的研究[J]. 中国机械工程, 2005, 16(14):1289-1291.
    \item 杨轶, 张宁欣, 任天令, 等. 集成铁电器件中的关键工艺研究[J]. 仪器仪表学报, 2003, 24(S4):192-193.
    \item Yang Y, Ren T L, Zhu Y P, et al. PMUTs for handwriting recognition. In press[J]. (已被Integrated Ferroelectrics录用)
  \end{achievements}


  \subsection*{专利}

  \begin{achievements}
    \item 任天令, 杨轶, 朱一平, 等. 硅基铁电微声学传感器畴极化区域控制和电极连接的方法: 中国, CN1602118A[P]. 2005-03-30.
    \item Ren T L, Yang Y, Zhu Y P, et al. Piezoelectric micro acoustic sensor based on ferroelectric materials: USA, No.11/215, 102[P]. (美国发明专利申请号.)
  \end{achievements}

\end{resume}


% 指导教师/指导小组评语
% 本科生不需要
% !TeX root = ../thuthesis-example.tex

\begin{comments}
% \begin{comments}[name = {指导小组评语}]
% \begin{comments}[name = {Comments from Thesis Supervisor}]
% \begin{comments}[name = {Comments from Thesis Supervision Committee}]

  论文提出了……

\end{comments}


% 答辩委员会决议书
% 本科生不需要
% !TeX root = ../thuthesis-example.tex

\begin{resolution}

  论文提出了……

  论文取得的主要创新性成果包括：

  1. ……

  2. ……

  3. ……

  论文工作表明作者在×××××具有×××××知识，具有××××能力，论文××××，答辩××××。

  答辩委员会表决，（×票/一致）同意通过论文答辩，并建议授予×××（姓名）×××（门类）学博士/硕士学位。

\end{resolution}


% 本科生的综合论文训练记录表（扫描版）
% \record{file=scan-record.pdf}

\end{document}
